\begin{textsample}{POS Dim 3 – to – Score 33.00 – to/to\_inf\_21.txt}  \label{ex:f3_pos_101}
Linguagem na Educação \textbf{Infantil} : experiências com a leitura e a escrita As \textbf{crianças}, no seu mundo de fantasias, experimentam, descobrem, planejam, formulam hipóteses, \textbf{testam}, \textbf{brincam}, interagem com seus \textbf{pares}, sejam eles, semelhantes ou diferentes, relacionam -se com os \textbf{adultos}, com a natureza, desenvolvem -se.
Nessa perspectiva, o trabalho pedagógico por meio de atividades que proporcionam vivências e experiências de leitura e de escrita na Educação \textbf{Infantil}, necessita compor -se na metáfora do texto de Malaguzzi ( 1999 ) para contemplar a essência de ser \textbf{criança}, de viver a \textbf{infância}, de garantir vez e voz a esses sujeitos e, ainda, de refletir acerca do papel da família, da escola e da sociedade nesse contexto.
As experiências das \textbf{crianças} com os elementos da sua cultura : modo de falar, de agir, de produzir bens materiais e imateriais, vão inserindo -as nas atividades de convívio social, quando, por exemplo, a \textbf{criança} imita a fala de um \textbf{adulto} em determinada situação real, por meio da \textbf{brincadeira} ou do faz de conta, quando pinta, \textbf{dramatiza} uma situação da vida cotidiana, a \textbf{criança} vai reproduzindo a sua cultura e acrescentando a ela novos significados, recriando a seu modo, sendo protagonista da sua língua, da sua fala e da sua produção cultural.
Arte de Petrus Habacuque Gualberto Fôlha – 04 anos Nesse contexto, e por meio das interações que estabelecem com outras \textbf{crianças}, elas se apropriam das diferentes formas de linguagens, observam seus próprios comportamentos e os comportamentos dos outros, diferem os papéis e as regras de convívio social nas mais diversas situações, e aprendem a respeitar a si mesmas e aos outros, a não ter preconceito, a argumentar a seu favor e a favor de outrem numa situação de disputa entre outras.
De acordo com as Diretrizes Curriculares Nacionais para Educação \textbf{Infantil} ( DCNEI ), as \textbf{instituições} devem garantir a todas as \textbf{crianças} vivências e experiências que promovam o conhecimento de si e do mundo por meio da ampliação de experiências \textbf{sensoriais}, expressivas, corporais que possibilitem \textbf{movimentação} ampla, expressão da individualidade e respeito pelos ritmos e \textbf{desejos} da \textbf{criança} ; a \textbf{imersão} nas diferentes linguagens e o progressivo domínio por elas de vários gêneros e formas de expressão : gestual, verbal, \textbf{plástica}, \textbf{dramática} e musical ; experiências de narrativas, de apreciação e interação com a linguagem oral e escrita, e convívio com diferentes suportes e gêneros textuais orais e escritos ( BRASIL, 2009, p. 27 ).
Ao praticar a leitura, quer sejam os códigos, imagens, cenas, paisagens, entre outras, a \textbf{criança} atribui sentido ao texto, e \textbf{consegue} relacioná -lo com o contexto e com as experiências que já vivenciou.
Ao participar das diversas vivências de leitura, a \textbf{criança}, enquanto ouvinte, inicia seu processo como futuro leitor.
Nessa perspectiva, toda leitura da palavra pressupõe uma leitura anterior do mundo, e toda leitura da palavra implica a \textbf{volta} sobre a leitura do mundo, de tal maneira que “ ler mundo ” e “ ler palavra ” constituam um movimento em que não há ruptura, em que você vai e \textbf{volta}.
E “ ler o mundo ” e “ ler palavra ”, no fundo, para mim, implicam reescrever o mundo.
Reescrever com aspas, quer dizer transformá -lo ( FREIRE ; BETTO 1986, p. 15 ).
Esse processo, no cotidiano escolar, pressupõe espaços e situações que garantam o contato das \textbf{crianças} com variados suportes e gêneros discursivos orais e escritos, com vistas a incentivar a exploração, o encantamento, a \textbf{curiosidade} e os questionamentos das \textbf{crianças} sobre a linguagem.
Para isto, é \textbf{interessante} que os espaços sejam convidativos e repletos de portadores de textos, obras de artes, produções com as \textbf{marcas} das \textbf{crianças}, com vistas a ampliar a \textbf{percepção} de mundo sobre a linguagem escrita.
Sendo a linguagem, um bem cultural que a \textbf{criança} começa a ter contato desde as suas primeiras experiências com o mundo exterior e a Educação \textbf{Infantil}, tendo a linguagem oral como a forma predominante verbal da \textbf{pequena} \textbf{infância}, como afirma Peter Moss ( 2009 ), não se pode furtar ao trabalho pedagógico que estimule a leitura, a contação de histórias, sejam elas reais ou \textbf{inventadas}.
A \textbf{creche} e a \textbf{pré-escola} vêm ocupando, nos últimos anos, o lugar de transmissão da cultura oral nas sociedades letradas, pois é nelas que os \textbf{adultos} têm o tempo e o espaço para sentar com as \textbf{crianças}, escutá -las e \textbf{conversar}.
É nela(s) que \textbf{adultos} e \textbf{crianças} sentam -se para ler e \textbf{ouvir} histórias, lendas, contos de fadas ; é lá também que circulam a cultura popular e a cultura \textbf{lúdica}, além de outros saberes que as \textbf{crianças} aprendem em suas culturas de \textbf{pares}, como jogos, canções, \textbf{brincadeiras}, e cantigas de \textbf{roda} que durante muitos séculos acompanharam o desenvolvimento humano ( BARBOSA ; DELGADO, 2012, p. 134 ).
Nessas interações e \textbf{brincadeiras} com a linguagem, as \textbf{crianças} vão percebendo as \textbf{marcas} da escrita, nos suportes que lhes são apresentados, nos suportes que observam, nos usos que a família e a escola fazem com o ler e o escrever e, naturalmente, iniciam o seu processo de ensaio da escrita, quando começam a desenhar letras em suas ilustrações.
É importante que os professores compreendam que o desenho \textbf{infantil} é a forma de escrita da \textbf{criança} e que cada faixa \textbf{etária} apresenta sua especificidade motora, cognitiva e cultural.
Nesse sentido, a \textbf{criança} representa simbolicamente o que vivencia em seu cotidiano, criando suas \textbf{marcas} e imprimindo significado ao longo de seu desenvolvimento.
Para tanto, é necessário compreender que a \textbf{criança} necessita de espaços apropriados para a produção gráfica, diferentes materiais riscantes e suportes que possibilitem a expressão de sua criatividade.
Todavia, o caminho até chegar a uma escrita autônoma é longo e complexo, portanto, considerando o processo de desenvolvimento da \textbf{criança} de 0 ( zero ) até 5 ( cinco ) anos de idade, mediante as transições que a \textbf{criança} vivencia, até ingressar no Ensino Fundamental, a pretensão da Educação \textbf{Infantil} deve ficar por conta das descobertas, como : o que se lê está registrado por diferentes sinais em algum suporte que podem representar diferentes \textbf{sons}, oportunizar diferentes emoções, que o seu próprio nome está escrito por estes sinais, e que estes sinais se \textbf{repetem} nos nomes das outras \textbf{crianças}, entre outros, tão importantes e carregados de sentidos para as \textbf{crianças}, que o processo vai além de decodificar códigos.
Nesse processo o que cabe então à Educação \textbf{Infantil}?
A ideia não é alfabetizar na perspectiva do Ensino Fundamental, conforme Barbosa ( 2017 ) “ a Educação \textbf{Infantil} deve trabalhar com práticas culturais que tenham a ver com a leitura, com a escrita, mas não é o momento de sistematizar a alfabetização ”.
Ou seja, a ideia não é forçar o processo de alfabetização das \textbf{crianças} com práticas pedagógicas que não respeitem as suas fases de desenvolvimento, mas de proporcionar às \textbf{crianças} o contato, a interação, o encantamento, a descoberta da leitura e da escrita, por meio de práticas pedagógicas que estimulem e garantam o direito das \textbf{crianças} de interagir com a cultura letrada e dela participar, nas diferentes situações cotidianas que envolvam o uso social e real da escrita, como por exemplo, a leitura coletiva dos bilhetes que vão para a casa ; a produção coletiva de um convite para um evento da turma ou da escola ; a leitura de cartazes enviados pelas unidades de saúde, entre outros.
A efetivação desse direito deve pautar -se nos eixos que orientam o trabalho com a \textbf{infância}, portanto, as interações e as \textbf{brincadeiras} são premissas para todas as atividades desenvolvidas pelas \textbf{crianças}, superando a visão instrucional e pragmática de trabalho com leitura e escrita, e entendendo que a apropriação da linguagem escrita requer elaboração cognitiva, reflexão sobre si e sobre o mundo, sobre conteúdos, estruturas, e as \textbf{crianças}, nessa faixa \textbf{etária}, ainda não desenvolveram a capacidade de isolar mentalmente um elemento ou uma propriedade de um todo, para considerar individualmente, ou seja, ainda não possuem capacidade de abstração.
Por isso, sugere -se trabalho com textos curtos, simples, com elementos que rimam, porém completos, considerando que a \textbf{criança} precisa de significação para compreender o todo.

% matched lemmas: adulto, brincadeira, brincar, conseguir, conversar, creche, criança, curiosidade, desejo, dramatizar, dramático, etário, imersão, infantil, infância, instituição, interessante, inventar, lúdico, marca, movimentação, ouvir, par, pequeno, percepção, plástico, pré-escola, repetir, roda, sensorial, som, testar, volta
\end{textsample}
